\documentclass[11pt,a4paper]{article}

\usepackage[margin=1in]{geometry}
\usepackage[T1]{fontenc}
\usepackage{textcomp}
\usepackage{amsmath,amssymb,amsthm}
\usepackage{mathtools}
\usepackage{booktabs}
\usepackage{array}
\usepackage{enumitem}
\usepackage{microtype}
\usepackage{xcolor}
\usepackage{hyperref}
\usepackage{listings}

\hypersetup{
  colorlinks=true,
  linkcolor=blue!65!black,
  citecolor=green!45!black,
  urlcolor=blue!75!black
}

\lstdefinelanguage{Lean}{
  keywords={structure, extends, def, noncomputable, theorem, lemma, inductive,
    where, let, in, match, with, fun, if, then, else, Prop, Type},
  sensitive=true
}

\lstset{
  language=Lean,
  basicstyle=\ttfamily\small,
  columns=fullflexible,
  breaklines=true,
  commentstyle=\color{gray!70!black},
  keywordstyle=\color{blue!70!black}\bfseries,
  stringstyle=\color{teal!70!black},
  frame=single,
  frameround=tttt,
  backgroundcolor=\color{gray!5}
}

\setlist{itemsep=3pt,topsep=5pt,parsep=0pt}

\newtheorem{theorem}{Theorem}[section]
\newtheorem{proposition}[theorem]{Proposition}
\theoremstyle{definition}
\newtheorem{definition}[theorem]{Definition}
\theoremstyle{remark}
\newtheorem{remark}[theorem]{Remark}

\newcommand{\HAomega}{\ensuremath{\mathrm{HA}^{\omega}}}
\newcommand{\SystemT}{\ensuremath{\mathrm{System\ T}}}
\newcommand{\N}{\ensuremath{\mathbb{N}}}
\newcommand{\Q}{\ensuremath{\mathbb{Q}}}
\newcommand{\Rep}{\ensuremath{\mathbf{Rep}}}
\newcommand{\mr}{\ensuremath{\;\mathbf{mr}\;}}
\newcommand{\Deriv}{\ensuremath{\mathrm{Deriv}}}
\newcommand{\Tm}{\ensuremath{\mathrm{Tm}}}
\newcommand{\extractClosed}{\ensuremath{\mathrm{extractClosed}}}
\newcommand{\lean}[1]{\texttt{\nolinkurl{#1}}}

\title{\textbf{From Existence Proofs to Executable Analysis:\\
Modified Realizability for \texorpdfstring{\HAomega}{HAomega} in Lean 4}}
\author{Ara Aslyan}
\date{\today}

\begin{document}

\maketitle

\begin{abstract}
We describe a Lean 4 development that turns constructive object-level proofs
into executable programs.  The formal core is a deep embedding of finite-type
Heyting arithmetic, with an intrinsically typed language of G\"odel System T
terms, a natural-deduction calculus, Kreisel modified realizability, a
certified extraction map, and code emitters for functional target languages.
The development is accompanied by exact rational numerical examples drawn
from Picard iteration and classical dynamical systems.  These examples are not
presented as new theorems of dynamics; their role is to show that the
proof-theoretic pipeline reaches running symbolic numerical code.

The central methodological point is axiom budgeting.  The extracted program
artifacts can be generated and evaluated without non-logical axioms in the
clean arithmetic examples, while the semantic soundness theorem and many
ambient analysis theorems have explicit Lean kernel footprints involving
\lean{propext}, \lean{Classical.choice}, and \lean{Quot.sound}.  For
constructive analysis this accounting becomes a calibration device: the
passage from continuous data with a uniform-continuity modulus to differentiable
data with explicit derivative moduli is measured by the footprint of
\lean{A0.toA1}.  The result is a proof-to-program framework in which both the
executable content of proofs and the represented data required by theorems are
made explicit.

For ordinary differential equations the pipeline produces solvers from proofs
whose statements name no algorithm. One derivation proves, for every closed
initial jet \(A\) and every polynomial--integral right-hand side \(F\), that
\(y = A + \int F(y)\) has a Taylor jet solving it below every degree; the proof
is an induction, and its extract is the Picard iteration. Instances include
\(\exp\), \(\tan\) (Riccati), \(y' = y^2\), and, for two-variable systems,
\(\sin\) and \(\cos\) and the Duffing oscillator; the extracts are proved to
compute the unique solution's coefficients. Every program in the catalog is
labelled by where its content comes from --- computed by the proof, an
explicit witness whose correctness is derived, or handed in by the statement.
\end{abstract}

\section{Introduction}

Constructive proof theory gives a precise form to a familiar slogan: an
existence proof can contain an algorithm.  Realizability and functional
interpretations make that slogan mathematical by assigning computational
witnesses to proofs of existential and universal statements
\cite{godel1958,kreisel1959,kohlenbach2008}.  Mechanized theorem proving adds
a further demand: the extraction procedure itself should be checked, not merely
trusted as an external program.

This article is written from a series of public showcase posts about a Lean 4
repository formalizing modified realizability for \HAomega.  Those posts
emphasize visible computations: exact rational Picard iterations, harmonic
oscillator approximations to sine and cosine, algebraic invariants for
Lotka--Volterra and Kepler systems, and dissipative or divergence identities
for Duffing, Van der Pol, and Lorenz systems.  The paper-level thesis is
slightly different and more precise:

\begin{quote}
A constructive derivation in a deeply embedded arithmetic can be translated
into a typed executable term, and the classical assumptions needed to justify
semantic correctness can be measured separately from the executable artifact.
\end{quote}

The development has three contributions.

\begin{enumerate}
\item It formalizes the syntax, typing, proof rules, modified-realizability
semantics, and extraction function for a finite-type arithmetic inside Lean 4.
\item It uses Lean's kernel axiom reporting as a discipline: extracted
programs, semantic soundness, and analysis bridges are not assigned one
blanket foundational status, but are measured declaration by declaration.
\item It organizes computable analysis by represented data.  The difference
between a continuous representation and a differentiable representation is not
the underlying mathematical function alone, but the moduli and witnesses that
the representation carries.
\item It extracts ODE solvers from proofs of algorithm-free specifications: a
single induction yields a solver for every polynomial--integral equation
\(y = A + \int F(y)\) and every two-variable system, with the extracted
programs proved to compute the unique solution, and with Cauchy, continuity and
differentiability moduli likewise computed by induction on the precision.
\end{enumerate}

The local build used while preparing this draft completed successfully at
7,923 jobs.  The exact count is not a mathematical contribution; it is included
only to identify the checked artifact state.

\section{The Formal Core}

The formal core separates object logic from Lean's ambient logic.  Terms,
formulas, derivations, realizers, and extracted programs are all represented as
Lean data.

\subsection{System T Terms}

The type language contains natural numbers, unit, products, and function
types:
\[
  \sigma,\tau ::= \mathrm{nat}\mid \mathrm{unit}\mid
  \sigma\times\tau\mid \sigma\to\tau .
\]
Terms are intrinsically typed: a term has type \lean{Tm Gamma tau}, indexed by
its object-language context and result type.  Thus a malformed object program
cannot be constructed.

\begin{lstlisting}[language=Lean]
inductive Tm : List Ty -> Ty -> Type where
  | var    : Var Gamma tau -> Tm Gamma tau
  | lam    : Tm (sigma :: Gamma) tau -> Tm Gamma (.arrow sigma tau)
  | app    : Tm Gamma (.arrow sigma tau) -> Tm Gamma sigma -> Tm Gamma tau
  | zero   : Tm Gamma .nat
  | succ   : Tm Gamma .nat -> Tm Gamma .nat
  | recNat : Tm Gamma tau ->
      Tm Gamma (.arrow .nat (.arrow tau tau)) ->
      Tm Gamma .nat -> Tm Gamma tau
  | pair   : Tm Gamma sigma -> Tm Gamma tau -> Tm Gamma (.prod sigma tau)
  | fst    : Tm Gamma (.prod sigma tau) -> Tm Gamma sigma
  | snd    : Tm Gamma (.prod sigma tau) -> Tm Gamma tau
  | star   : Tm Gamma .unit
\end{lstlisting}

The actual repository extends this core with rational arithmetic and additional
term constructors used by the numerical examples.

\subsection{Object-Level Derivations}

\lean{Deriv} is the natural-deduction calculus for the embedded \HAomega.
Rules include implication, conjunction, universal and existential quantifiers,
equality reasoning, and induction over natural numbers.  A derivation is not a
Lean theorem directly proving an ambient proposition; it is a proof object in
the formalized object logic.

This distinction matters.  It lets the development inspect the proof as data,
translate it into a realizer, and prove a separate soundness theorem connecting
object-level derivability to Lean-level semantic truth.

\subsection{Modified Realizability}

Kreisel modified realizability assigns each formula a realizer type and a
predicate describing when a value realizes the formula.  Computational content
appears exactly where the logic demands it: a proof of an existential statement
contains a witness, and a proof of a universal statement contains a function.

\begin{lstlisting}[language=Lean]
def MR : (phi : Formula Gamma a) ->
    Env Gamma -> (realizerTy phi).interp -> Prop
  | .eq s t,   e, _ => s.eval e = t.eval e
  | .and p q,  e, r => MR p e r.1 /\ MR q e r.2
  | .imp p q,  e, r => forall x, MR p e x -> MR q e (r x)
  | .all sigma p, e, r =>
      forall x : sigma.interp, MR p (Env.cons x e) (r x)
  | .ex sigma p, e, r => MR p (Env.cons r.1 e) r.2
\end{lstlisting}

\section{Extraction and the Smallest Complete Example}

The extraction map sends object-level derivations to System T terms:
\[
  \extractClosed(D) : \Tm\ []\ (\operatorname{realizerTy}(\varphi))
\]
for a closed derivation \(D : \Deriv\ []\ \varphi\).

The cleanest reproducible example is the declaration
\lean{doublingRealizer} in \lean{HAomega/AnalysisDeriv.lean}.  The derivation
\lean{doublingRecDeriv} proves the totality of an iteration sequence with
initial value \(1\) and step \(x\mapsto x+x\).  Extraction produces a closed
System T program whose witness component computes \(2^n\).

\begin{lstlisting}[language=Lean]
def doublingStepTm : Tm [] (.arrow .nat .nat) :=
  .lam (.add (.var .here) (.var .here))

def doublingRecDeriv :
    Deriv .nil (.all .nat (iterInv .nat [] (.succ .zero) doublingStepTm)) :=
  iterSequenceD .nat (.succ .zero) doublingStepTm

def doublingRealizer : Tm [] (.arrow .nat (.prod .nat .unit)) :=
  extractClosed doublingRecDeriv

def doublingExtracted (n : Nat) : Nat :=
  ((extractClosed doublingRecDeriv).eval Env.nil n).1
\end{lstlisting}

The Lean file checks the following values by kernel reduction:
\[
0\mapsto 1,\quad
1\mapsto 2,\quad
2\mapsto 4,\quad
3\mapsto 8,\quad
4\mapsto 16,\quad
5\mapsto 32,\quad
10\mapsto 1024 .
\]

This example is deliberately modest.  Its value is provenance: the derivation,
extraction, evaluation, and axiom report all refer to named declarations in
the repository.

\section{Axiom Footprints}

The development should not be summarized by the sentence ``everything is
axiom-free.''  That is false for many ambient analysis theorems.  The correct
claim is a separation.

\begin{table}[htbp]
\centering
\begin{tabular}{p{0.26\linewidth}p{0.25\linewidth}p{0.40\linewidth}}
\toprule
\textbf{Declaration or layer} & \textbf{Role} & \textbf{Measured footprint} \\
\midrule
\lean{extractClosed} & syntactic extraction & no non-logical axioms \\
\lean{doublingRealizer} & extracted arithmetic realizer & no non-logical axioms \\
\lean{soundnessClosed} & semantic soundness & \lean{propext}, \lean{Classical.choice}, \lean{Quot.sound} \\
\lean{eftc1}, \lean{eftc2_thm} & analysis theorems & \lean{propext}, \lean{Classical.choice}, \lean{Quot.sound} \\
dynamical-system identities & algebraic/rational examples & generally \lean{propext}, \lean{Classical.choice}, \lean{Quot.sound} \\
\bottomrule
\end{tabular}
\caption{The useful claim is not a blanket zero-axiom slogan, but a measured
separation between executable artifacts and semantic or analytic bridges.}
\end{table}

This axiom accounting is one of the main methodological contributions.  It
makes the computational artifact and the proof of its semantic adequacy
separate objects of study.

\section{Representations as Data Budgets}

The analysis layer asks a calibration question:

\begin{quote}
How much represented data must a theorem carry before its constructive proof
yields executable numerical content?
\end{quote}

The representation \(A_0\) carries an exact rational evaluator together with a
modulus of uniform continuity.  The representation \(A_1\) carries additional
derivative data and a modulus of differentiability.  The first fundamental
theorem, as formalized by \lean{eftc1}, is available at the \(A_0\) level:
integration needs the continuity modulus.  The second fundamental theorem, as
formalized by \lean{eftc2_thm}, is available at the \(A_1\) level: it requires
the additional differentiability data.

The bridge
\begin{lstlisting}[language=Lean]
noncomputable def A0.toA1 (A : A0) (h : HasUnifDeriv A) : A1
\end{lstlisting}
uses classical choice to select missing witnesses.  This is the precise place
where the informal slogan ``the modulus is the content'' becomes formal.  The
mathematical function is not enough; the representation must carry the
quantitative data needed by the theorem.

\begin{remark}
The safe formulation is that the axiom footprint of \lean{A0.toA1} exhibits
the gap identified by Myhill's computability result \cite{myhill1971}.  It
does not constitute an internal formal proof of Myhill's negative theorem,
because the present object theory does not include a computability predicate
able to state that negative result.
\end{remark}

\section{Galois Adequacy}

The repository also contains an algebraic layer of represented spaces.  A
representation packages a carrier of codes, an approximation relation, and an
equivalence relation on codes.  Morphisms preserve approximation with an
explicit precision transformer.  A retract preorder \(R_1 \preceq R_2\)
organizes when one representation can be used through another.

The bridge theorem \lean{central_adequacy_theorem} connects the
proof-theoretic and representation-theoretic layers.  Given a closed \HAomega
derivation of a \(\forall\exists\)-formula and a local compatibility
specification \lean{RepAdequacySpec}, the theorem packages the extracted code
as a representation-level adequate realizer.

This is the point at which the work differs from ordinary examples of
symbolic computation.  The emitted code is downstream of a formal derivation,
and the representation layer records how much precision must be requested from
input codes to produce an output approximation.

\section{Executable Analysis Case Studies}

The showcase posts are best read as examples of the executable end of the
pipeline.

\paragraph{Harmonic oscillator.}
The harmonic oscillator
\[
  y''+y=0
\]
is represented as a first-order system and iterated by an exact rational
Picard operator over polynomial coefficient lists.  Lean checks, among other
values, the approximations
\[
  \cos(1/2) \approx 337/384,\qquad
  \sin(1/2) \approx 1841/3840 .
\]
The formal claim is exact rational Picard computation and checked arithmetic,
not a new theorem about trigonometric functions.

\paragraph{Van der Pol.}
The file \lean{VanDerPol.lean} proves the divergence/trace identity
\lean{vanderpol_divergence_trace} and checks rational Picard trajectory
values.  The global limit-cycle theorem belongs to the mathematical
background, not to the current formal claim.

\paragraph{Lotka--Volterra.}
The file \lean{LotkaVolterra.lean} proves the algebraic cancellation
\lean{lotka_volterra_invariant_cancel} and checks initial rational
trajectory responses.  The invariant gives the example its mathematical
meaning; the formal artifact is the cancellation identity plus kernel-checked
computation.

\paragraph{Duffing.}
The file \lean{Duffing.lean} proves \lean{duffing_energy_dissipation_id},
expressing the dissipative energy identity, and checks rational Picard values.
The current formal scope is the identity and the computations, not a full
basin-of-attraction analysis.

\paragraph{Lorenz.}
The file \lean{Lorenz.lean} proves \lean{lorenz_volume_contraction_rate}.  A
paper should not say that Lean proves the existence of a strange attractor
here; it proves the divergence identity and computes exact initial Picard
segments.

\paragraph{Kepler.}
The file \lean{Kepler.lean} proves \lean{kepler_angular_momentum_conserved}.
The separate symplectic integrator checks angular momentum preservation for
selected discrete trajectories.  The formal contribution is conservation plus
exact rational computation.

\subsection{Differential-Integral Jet Algebras and the Classification of Their Realizers}

Beyond finite-degree polynomial stepping, the modules \lean{ClosedCalculus.lean},
\lean{HolonomicCS.lean}, \lean{MajorantDom.lean}, \lean{JetCauchyDeriv.lean},
and \lean{ClosedFourierODE.lean} formalize formal power series jets
\(\lean{JetTy} = \mathrm{nat}\to\mathrm{rat}\) inside \SystemT\ together with
two Cauchy-majorized value-level algebras:
\begin{itemize}
\item \lean{AnalyticJet} (\(|c_n|\cdot n! \le 2^M(2^R)^n\), entire functions of
exponential type on \([-1,1]\)), and
\item \lean{RadiusJet} (\(|c_n| \le 2^M(2^R)^n\), local Cauchy--Hadamard germs of
radius \(\rho = 2^{-R}\), closed under addition, Cauchy multiplication,
differentiation, integration, \(1/(1-G)\) inversion, composition, and nonlinear
Picard fixed points).
\end{itemize}
Both algebras embed into the approximating differentiable representation
\(\lean{E1}\) (\lean{AnalyticJet.toE1}, \lean{RadiusJet.toE1}), where
\(\lean{E1.diff}\) quantifies over an approximating derivative family
\(\exists F : \N\to\Q\to\Q\) rather than an exact rational function.

Across the thirty-seven extracted programs cataloged in \lean{ShowAll.lean} and
\lean{EXTRACTED_HAOMEGA.md}, we distinguish three classes of object-level
derivations, by where the computational content comes from:
\begin{enumerate}
\item \textbf{Proof-computed realizers}, where the specification names no
algorithm and the program's recursion is assembled by extraction from
induction, \lean{tiEps0}, or case analysis. These are the discrete case
studies (gcd, Goodstein, Kirby--Paris, Sperner, Hercules) and the
square-root approximation, and one jet program: \lean{picardExpSolveD}
(\lean{HAomega/PicardExpSolve.lean}) proves
\(\forall k\,\exists Y.\,\forall i < k.\; \mathcal{P}(Y)(i) = Y(i)\) for
\(y' = y\), \(\mathcal{P}(Y) = 1 + \mathcal{I}Y\), a specification that names the
integral equation but no algorithm. The proof is \lean{Deriv.ind} on \(k\),
building the witness at \(k+1\) as \(\mathcal{P}\) of the witness at \(k\);
the extracted solver is the Picard recursion, and equals the iterate.
The same holds for the nonlinear \(y' = y^2\) (\lean{quadBlowupSolveD},
\(\mathcal{P}(Y) = 1 + \mathcal{I}(Y\cdot Y)\)), where the induction rests on a
prefix lemma derived inside the calculus: a Cauchy-product coefficient depends
only on the factors' coefficients up to the same index. The Riccati equation
\(y' = 1 + y^2\) (\lean{riccatiSolveD}) and the Duffing system
\(u' = v,\ v' = -u - u^2\) (\lean{duffingSolveD}, on pairs of jets) follow by
the same argument.
All the scalar cases are instances of one derivation: for every closed initial
jet \(A\) and every right-hand side \(F\) built from \(Y\), constants,
\(+\), \(-\), \(\times\) and \(\mathcal{I}\), \lean{genSolveD} proves
\(\forall k\,\exists Y.\,\forall i < k.\; (A + \mathcal{I}F(Y))(i) = Y(i)\)
by induction, resting on a prefix lemma (\lean{JExpr.prefixT}) derived by
recursion on \(F\); its extract is the Picard iteration. The same holds
for systems in two jet variables (\lean{genSolve2D}); sine and cosine arise as
its instance \(u' = v,\ v' = -u\), the Duffing oscillator as another. Finally
\lean{genSolveND} handles systems of any size: a list of equations
\((A_v, F_v)\), with equation \(v\) selected inside \SystemT\ by recursion on
\(v\).
\item \textbf{Explicit witness, derived correctness}, where the specification
names no algorithm, but the derivation writes the witness at \lean{Deriv.exI}
and the subsequent induction or scaling rules prove contentless equations, so
they certify the witness without contributing computation:
\begin{itemize}
\item \lean{jetCauchyModulusGenuineD} (\lean{HAomega/JetCauchyDeriv.lean})
writes the Cauchy modulus \(N(k) = k + (B + 3)\) and justifies it from an
\(N\)-free, per-term majorant hypothesis
\(\forall n,\; |S_n(x)-S_{n+1}(x)| < 2^{B+3}(2^{-(n+1)}-2^{-(n+2)})\),
deriving the multi-step tail bound by object-level induction (\lean{Deriv.ind})
on \(d_2\), telescoping cancellation, and precision weakening. Four closed
\(\mathbb{Q}\)-laws (zero bound, triangle inequality, telescoping, scale
cancellation) enter as a hypothesis proved in the metatheory. The majorant
hypothesis is discharged for both \lean{AnalyticJet} and \lean{RadiusJet};
\item \lean{jetContModulusD} and \lean{jetDiffModulusD} write the continuity
modulus \(\omega(k) = k + (B_1 + 1)\) and differentiability modulus
\(\delta(k) = k + (B_2 + 1)\) and certify them via \lean{Deriv.convQLipScale},
discharged for both \lean{AnalyticJet} and \lean{RadiusJet}.
\end{itemize}
\item \textbf{Witness-style realizers}, where the theorem statement names the
target term and the derivation supplies it to \lean{Deriv.exI}:
\lean{ftcAntiderivD} (which hands \lean{tmJetInteg} to \lean{exI} and derives
termwise cancellation \(\mathcal{D}(\mathcal{I}c)(n) - c(n) = 0\) via
\lean{Deriv.convQMulDivSubZero}), \lean{jetMulD}, \lean{jetEvalPolyD}, the
quadratic Newton--Schulz jet inverter \lean{newtonJetInvD} (linked to
\lean{RadiusJet.invOneSub} by \lean{extractedNewtonJetInv_agree}), and the
Picard iterations \lean{picardJetIterD}, \lean{riccatiPicardD},
\lean{quadBlowupPicardD}, \lean{duffingPicardD}, which prove
\(\forall k\,\exists Y.\; Y = \mathcal{P}^k(Y_0)\) by supplying the iterate.
Degree locking --- that \(k\) Picard steps fix \(k\) Taylor coefficients --- is
proved in the metatheory for a generic guarded step
(\lean{picard_jetAgree_lock}) and instantiated at all four extracted solvers
(\lean{HAomega/PicardLocking.lean}): iterate \(k+1\) agrees with iterate \(k\)
below degree \(k\). These are theorems about the extracted programs, not part
of the derivations.
\end{enumerate}

\subsection{Toward a Closed Universe}

The longer-range goal is a minimal universe inside \HAomega\ in which
\(+, -, \times, \frac{d}{dx}, \int\), \(\exp\), \(\sin\), \(\cos\), Taylor
series and Fourier series live together, with existence statements yielding
extracted programs. The general solvers above are its generating mechanism:
the universe is the closure of constants under \(+,-,\times,\mathcal{I}\) and
under solving \(y = A + \mathcal{I}F(y)\). The following are mechanized.
\begin{itemize}
\item \textbf{Uniqueness.} Any solution of \(Y = A + \mathcal{I}F(Y)\) agrees
with the \(k\)-th extracted iterate below degree \(k\)
(\lean{genSolve_unique_val}, \lean{genSolve2_unique_val}); the extracts compute
\emph{the} Taylor coefficients, not merely those of some solution.
\item \textbf{Inverses.} On the unnormalized \(\mathbb{Q}\) the guard
\(G(0)\cdot Y(i)\) of \(Y = 1 + G\cdot Y\) is not structurally zero, so the guard
is taken to be the shift \(x\cdot F\), which needs no division.
\lean{shiftSolveD} proves \(\forall k\,\exists Y.\,\forall i<k.\;
(A + x\cdot F(Y))(i) = Y(i)\) for every \(A, F\), giving \(1/(1 - xH)\) for every
\(H\); for \(1/(1 - x - x^2)\) the extract yields the Fibonacci numbers exactly.
\item \textbf{A computed Cauchy modulus.} The \(\forall k\,\exists N\) of
the jet Cauchy theorem is also proved by induction on the precision \(k\)
(\lean{jetCauchyModulusComputedD}): if the tail bound from \(N\) is below
\(2^{-k}\), advancing \(N\) by one halves it. The extracted modulus is the
recursion this induction produces. The continuity and differentiability
moduli are computed the same way (\lean{jetContModulusComputedD},
\lean{jetDiffModulusComputedD}), without further rational laws.
\item \textbf{Composition.} \(f\circ(G - G(0))\) for an ODE-defined \(f\) solves
\(Y = f(0) + \mathcal{I}(G'\cdot F(Y))\), an instance of \lean{genSolveD} with the
constant jet \(G'\). For \(f = \exp\) the extract is proved equal to the
metatheory's formal composition for every \(G\), by the chain rule
\(\frac{d}{dx}\exp(G - G(0)) = G'\exp(G - G(0))\) and uniqueness
(\lean{expComp_eq_compZero}).
\item \textbf{The Fourier layer.} Trigonometric polynomials
\(\sum_{|k|\le N} c_k e^{ikx}\) over \(\mathbb{Q}[i]\) are exactly closed under
\(+\), \(\times\), \(\frac{d}{dx}\) and the mean-zero \(\int\); infinite Fourier
series are not, since a product coefficient is an infinite sum. A Taylor bridge
is proved to commute with \(\frac{d}{dx}\) and (for real polynomials) with
\(\times\), and to send the modes \(\cos(kx)\), \(\sin(kx)\) to their Taylor
jets. The modes themselves are proof-computed, as the instance
\(u' = m v,\ v' = -m u\) of \lean{genSolve2D}, and agree with the bridge
(\lean{genSolve2_mode_trig}).

\item \textbf{The algebra inside the object logic.} Structural equality on the
unnormalized rationals cannot state \(x + 0 = x\), so the laws are stated with a
value equality \(s \approx t\), \emph{defined} as \(s + 0 = t + 0\); its rules
(every operation returns a normal form and reads only the values of its
arguments; the field laws a bare variable blocks) enter by one schema rule,
\lean{Deriv.qlaw}, sound by uniqueness of normal forms. Derived in the calculus:
\(\exp' \approx \exp\), \(\sin' \approx \cos\), \(\cos' \approx -\sin\),
FTC1 \(D(\mathcal{I}c) \approx c\) and FTC2 \(\mathcal{I}(Dc)(n+1) \approx c(n+1)\)
for every jet \(c\) (\lean{QEquiv.lean}), and the product rule
\(D(fg) \approx Df\cdot g + f\cdot Dg\) for every \(f, g\), by induction over the
Cauchy sum (\lean{productRuleLawD}).
\item \textbf{Bisection and roots.} \(\forall j\,\forall c^{\N\to\N}\,\forall lo.\;
c(lo) \ne c(lo + 2^j) \to \exists i < 2^j.\; c(lo+i) \ne c(lo+i+1)\) is proved by
induction on \(j\) (\lean{bisectD}); its extract is binary search, a type-2
operator certified through soundness. A first proof of the same statement,
with the induction hypothesis used inside both branches of the case split,
extracts an exponential program; the cost is a property of the proof's shape
under extraction. Composed with the certified evaluator, the operator gives
approximate roots of every expression of the universe on its certified disk
(\lean{UExpr.root_spec}), e.g.\ \(\ln(41/40)\) and \(\arcsin(1/5)\).
\end{itemize}
Not yet inside the object logic: the chain rule and \(\sin^2+\cos^2=1\).

\section{Related Work}

This work is adjacent to several established lines.  Minlog provides program
extraction from constructive proofs, including examples in exact real
computation, but its extraction implementation is external to Lean's kernel
\cite{schwichtenberg2012,berger1998}.  C-CoRN formalizes constructive analysis
inside Coq's ambient type theory \cite{cruzfilipe2004}; the present work
instead emphasizes a deep object logic and an internal extraction theorem.
Kohlenbach's proof mining program and Cheval's Lean formalization of
Dialectica-style proof mining are close in proof-theoretic spirit
\cite{kohlenbach2008,cheval2023}; the present development uses modified
realizability and tracks extracted terms directly.

Represented-space approaches such as Incone \cite{steinberg2017,steinberg2022}
and Weihrauch reducibility \cite{weihrauch2000,brattka2021} are closest to the
analysis and representation layer.  The intended distinction is that this
repository combines represented spaces with a deep \HAomega proof calculus and
an extraction-to-realizer bridge, while Weihrauch theory primarily orders
computational problems rather than the represented data budget itself.

Treating a function through its stream of Taylor coefficients, with
differentiation and integration as stream operations, is classical:
Pavlovi\'c and Escard\'o give a coinductive formalization of elementary
calculus covering Taylor series and ordinary differential equations
\cite{pavlovic1998}, and McIlroy implements power-series algorithms, including
solutions of differential equations, as lazy functional programs
\cite{mcilroy1999}. Our jets are such streams; what is added here is that the
solvers are extracted from object-level proofs and certified. Polynomial
ordinary differential equations are also a model of computation: Bournez,
Gra\c{c}a and Pouly characterize polynomial time by solutions of polynomial
ODEs of polynomial length \cite{bournez2017}. The class of right-hand sides our
general solver accepts is polynomial in the unknown jets, which places the
extracted solvers in that setting; we make no complexity claim.

\section{Limitations}

Several boundaries should be explicit.

\begin{itemize}
\item The development does not prove PA-independence of Goodstein or
Kirby--Paris style statements.
\item It does not internally prove Myhill's negative theorem; it exhibits the
relevant gap as an axiom dependency of a bridge construction.
\item The dynamical-system case studies do not currently prove global limit
cycles, strange attractors, or basin convergence.
\item Picard and numerical terms should not be called extracted unless the
specific artifact is the result of \lean{extractClosed} applied to a
\lean{Deriv}.
\item Build counts are audit metadata, not evidence of mathematical novelty.
\item Equality at type \(\mathrm{rat}\) in the object logic is structural on an
unnormalized representation, so value-true identities can be false as stated
(\(x + 0 = x\) fails). The ODE results are therefore stated in integral form,
inverses use a shift guard, and the algebra laws of the closed universe are not
yet stated inside the logic; whether to normalize \(\mathbb{Q}\) or add a value
equality is an open design decision.
\item The specifications of the proof-computed solvers name the equation (the
Picard operator), though not the iteration; the differential form
\(Y' = F(Y)\) is not derived inside the logic.
\item Bridge, uniqueness and Fourier theorems are metatheory theorems about
values and use \lean{Classical.choice} through Mathlib's \(\mathbb{Q}\); the
extracted programs themselves stay \lean{[propext, Quot.sound]}.
\end{itemize}

\section{Conclusion}

The project shows that constructive proof extraction, exact symbolic numerical
computation, code emission, and axiom-footprint accounting can be maintained in
one checked Lean artifact.  The most durable contribution is not any one
numerical example.  It is the discipline of separating proof syntax, extracted
programs, semantic soundness, and represented data, so that the computational
content of a theorem can be inspected rather than inferred.

\bibliographystyle{plain}
\bibliography{references}

\end{document}
